% ============================================================
%  Chapter 7: Mass as Accumulated Non-Actualisation
%  Source: bounded-phase-space-trajectory.tex
%  Consequences 19--25 and surrounding material
% ============================================================

\epigraph{%
  Mass is not a primitive quantity.\\
  It is the normalised rate at which a persistent bounded oscillator\\
  fails to actualise all of its categorical possibilities\\
  per unit oscillation time.%
}{}

With the Lorentz group and the dispersion relation established in
Chapter~\ref{ch:lorentz}, the path to mass is short.  Section~\ref{sec:restmass}
reads off the rest mass from the dispersion relation and derives the identity
$m_{0} = \hbar\omega_{0}/c^{2}$.  Section~\ref{sec:loschmidt} pivots to the
thermodynamic consequences of the accumulative structure of partition
histories, proving the Loschmidt principle (entropy is non-decreasing) and
the arrow of time as a set-theoretic theorem.
Section~\ref{sec:massmemory} provides the deepest interpretation: mass is the
normalised rate of categorical non-actualisation per oscillation cycle.
Sections~\ref{sec:emc2} and~\ref{sec:equivmass} derive mass--energy equivalence
and the equivalence of inertial and gravitational mass.

\section{Inertial Mass from the Dispersion Relation}
\label{sec:restmass}

\begin{consequence}{Rest Mass Identity}{restmass}
The rest mass of a Koopman eigenmode with rest-frame angular frequency
$\omega_{0}$ is
\begin{equation}
  m_{0} = \frac{\hbar\omega_{0}}{c^{2}}.
  \label{eq:restmass}
\end{equation}
\end{consequence}

\begin{proof}
The group velocity of the mode is obtained from the dispersion relation
$\omega^{2} - c^{2}k^{2} = \omega_{0}^{2}$ (Consequence~\ref{cons:dispersion}):
\[
  v = \frac{d\omega}{dk} = \frac{c^{2}k}{\omega}.
\]
Define the four-momentum components $E := \hbar\omega$ and $p := \hbar k$.
Substituting,
\[
  v = \frac{c^{2}p}{E},
\]
equivalent to the relativistic mechanical relation $p = Ev/c^{2}$.  In the
non-relativistic limit $v \ll c$, Taylor-expand the dispersion relation:
\[
  E = \hbar\omega
    = \hbar\sqrt{\omega_{0}^{2} + c^{2}k^{2}}
    = \hbar\omega_{0}\sqrt{1 + \frac{c^{2}k^{2}}{\omega_{0}^{2}}}
    \approx \hbar\omega_{0} + \frac{\hbar c^{2}k^{2}}{2\omega_{0}}
    = \hbar\omega_{0} + \frac{p^{2}v}{2p}
    = \hbar\omega_{0} + \tfrac{1}{2}\frac{p^{2}m_{0}}{\hbar\omega_{0}/c^{2}}.
\]
Identifying the $v$-independent part $\hbar\omega_{0}$ with the rest energy
$m_{0}c^{2}$ (anticipating Consequence~\ref{cons:emc2}, which confirms this
identification), we obtain $m_{0}c^{2} = \hbar\omega_{0}$, i.e.,
$m_{0} = \hbar\omega_{0}/c^{2}$.
\end{proof}

\subsection{Numerical confirmation}

Consequence~\ref{cons:restmass} asserts $m_{0} = \hbar\omega_{0}/c^{2}$.
The predictive content is that a single angular frequency $\omega_{0}$
suffices to predict the mass: the ratio $m_{0}/\omega_{0} = \hbar/c^{2}$
is universal.

\begin{center}
\begin{tabular}{l|c|c|c}
\toprule
Object & $\omega_{0}$ (rad/s) & $m_{0} = \hbar\omega_{0}/c^{2}$ (kg) & Measured (kg) \\
\midrule
Electron & $7.76\times 10^{20}$ & $9.11\times 10^{-31}$ & $9.109\times 10^{-31}$ \\
Proton   & $1.42\times 10^{24}$ & $1.67\times 10^{-27}$ & $1.6726\times 10^{-27}$ \\
Muon     & $1.61\times 10^{23}$ & $1.88\times 10^{-28}$ & $1.884\times 10^{-28}$ \\
\bottomrule
\end{tabular}
\end{center}

In each row $\omega_{0}$ is back-calculated from the measured rest mass via
$\omega_{0} = m_{0}c^{2}/\hbar$.  The universal ratio is
$\hbar/c^{2} = 1.173\times 10^{-51}$~kg\,s, confirmed in every row.

\subsection{Why mass is positive}

Consequence~\ref{cons:restmass} gives $m_{0} = \hbar\omega_{0}/c^{2}$ with
$\omega_{0} > 0$ (Consequence~\ref{cons:restfreq}); hence $m_{0} > 0$.  The
positivity of mass is not a postulate; it is a consequence of the rest
frequency being positive, which in turn is a consequence of localisation in a
bounded region.

A would-be negative-mass object would require $\omega_{0} < 0$, contradicting
the positivity of oscillation frequencies in Koopman spectral theory
(frequencies are real and non-negative by convention).  Within the present
framework, negative mass is inadmissible without a generalisation of the Axiom.

\section{The Loschmidt Principle and the Arrow of Time}
\label{sec:loschmidt}

Before deriving the interpretive content of mass, we establish the
thermodynamic direction that mass accumulation presupposes.

\begin{definition}[Partition structure of a state]
\label{def:partstruct}
For a state $S$ at time $t$, the \emph{partition structure}
$\Se(S, t)$ is the set of partition cells occupied by $S$ at all depths $n$
of the hierarchy of Axiom~\ref{ax:bpsl}(iii).
\end{definition}

\begin{consequence}{Loschmidt Principle}{loschmidt}
Let $S_{1}$ be a state at time $t_{1}$ and $S_{2}$ the same state at time
$t_{2} > t_{1}$, evolved under the measure-preserving flow $\phi_{t}$.  Then
\begin{equation}
  \Se(S_{1}, t_{1}) \subseteq \Se(S_{2}, t_{2}),
  \label{eq:loschmidt}
\end{equation}
with strict inclusion on a set of positive measure.  Entropy
$S_{\mathrm{part}}(t) := \kB\ln|\Se(S, t)|$ is non-decreasing and strictly
increasing on a set of positive measure.
\end{consequence}

\begin{proof}
At $t_{1}$, $S$ has visited precisely the cells in $\Se(S_{1}, t_{1})$.  At
$t_{2} > t_{1}$, $S$ has visited the same cells plus any additional cells
traversed between $t_{1}$ and $t_{2}$.  Hence
$\Se(S_{1}, t_{1}) \subseteq \Se(S_{2}, t_{2})$.

\emph{Strict inclusion.}  By Consequence~\ref{cons:oscillatory}, the
trajectory visits arbitrary neighbourhoods of its starting point, and by the
elimination of static dynamics (Lemma~\ref{lem:staticelim}) the trajectory
is not constant.  Hence on a time interval of positive measure the trajectory
visits cells not in $\Se(S_{1}, t_{1})$.  These cells are added to
$\Se(S_{2}, t_{2})$ and cannot be removed: a measure-preserving flow does
not delete the past record.  Formally, ``reversing'' the accumulation would
require deleting partition cells from $\Se$, which is forbidden by
Axiom~\ref{ax:bpsl}(iii) (partitions are refined, not coarsened, under the
system's observed history).

The entropy $S_{\mathrm{part}}(t) := \kB\ln|\Se(S, t)|$ is therefore
non-decreasing, with strict increase on sets of positive measure.
\end{proof}

\begin{remark}[Identity from non-actualisation]
The complement of $\Se(S, t)$ in the full partition algebra is the set of
cells $S$ has \emph{not} visited.  This complement is the negative-space
identity of $S$: what distinguishes $S$ from other states is what it is
\emph{not}, and the measure of this ``not'' set shrinks monotonically with
time.  Entropy increase is therefore a logical necessity of the monotonicity
of accumulation, not a probabilistic tendency.  The Boltzmann H-theorem
\citep{boltzmann1877} is a statistical corollary of
Consequence~\ref{cons:loschmidt}.
\end{remark}

\begin{corollary}[Second law of thermodynamics]
\label{cor:secondlaw}
The total entropy of an isolated bounded system is a non-decreasing function
of time.
\end{corollary}

\begin{proof}
Immediate from Consequence~\ref{cons:loschmidt}:
$\Se(S, t_{1}) \subseteq \Se(S, t_{2})$ implies
$|\Se(S, t_{1})| \le |\Se(S, t_{2})|$, hence
$S_{\mathrm{part}}(t_{1}) \le S_{\mathrm{part}}(t_{2})$.
\end{proof}

\begin{remark}
The second law of thermodynamics is usually stated as an empirical observation
without derivation.  In the present framework it is a Consequence of the Axiom
via the Loschmidt principle.  No additional assumption of microscopic
reversibility, ergodicity, or probability is required; the monotonicity is a
purely combinatorial fact about accumulated partition structure.
\end{remark}

\begin{proposition}[Arrow of time]
\label{prop:arrowtime}
The direction in which entropy increases defines the arrow of time.
\end{proposition}

\begin{proof}
Consequence~\ref{cons:loschmidt} asserts
$\Se(S, t_{1}) \subseteq \Se(S, t_{2})$ for $t_{1} < t_{2}$.  The reverse
inclusion, $\Se(S, t_{2}) \subseteq \Se(S, t_{1})$, would require deletion of
cells from the accumulated history.  Such deletion violates
Axiom~\ref{ax:bpsl}(iii), which refines (not coarsens) the partition.  Hence
the asymmetry $t_{1} < t_{2} \Rightarrow \Se(S, t_{1}) \subseteq \Se(S, t_{2})$
defines a uniquely oriented temporal ordering.
\end{proof}

\subsection{The third law of thermodynamics}

\begin{consequence}{Third Law}{thirdlaw}
As $T \to 0$, the categorical entropy of a system with a unique ground state
(a single most-refined partition cell) approaches zero:
$\lim_{T\to 0} S_{\mathrm{cat}} = 0$.
\end{consequence}

\begin{proof}
At $T = 0$, the system occupies its lowest-energy partition cell with no
thermal fluctuations.  The accumulated partition history $\Se$ consists of a
single cell (the ground state), so
$S_{\mathrm{cat}} = \kB \cdot 1 \cdot \ln 1 = 0$ (with $n = 1$ accessible
cell).  The third law $S \to 0$ at $T = 0$ is a Consequence of the Axiom
whenever the ground state is unique.
\end{proof}

\begin{remark}
Nernst's theorem states that the entropy of a perfect crystal approaches zero
as temperature approaches absolute zero.  Consequence~\ref{cons:thirdlaw} is
the same statement framed categorically; it is a Consequence of the Axiom
rather than an independent postulate.
\end{remark}

\section{Mass as Accumulated Non-Actualisation}
\label{sec:massmemory}

\begin{consequence}{Mass as Accumulated Non-Actualisation}{massmemory}
In a system oscillating at rest frequency $\omega_{0}$, each cycle partitions
the local state space into $b^{d}$ categorical cells (where $b$ is the
branching base and $d$ the number of partitionable degrees of freedom), of
which one is actualised and $b^{d} - 1$ are not.  The accumulated
non-actualisation rate
\begin{equation}
  \rho_{\mathrm{na}} := (b^{d} - 1)\,\frac{\omega_{0}}{2\pi}
  \quad\text{cells per second}
  \label{eq:rateNA}
\end{equation}
has the dimensional signature of an inverse relaxation time, and its
normalisation by $c^{2}/\hbar$ produces the rest mass:
\begin{equation}
  m_{0}
  = \frac{\hbar}{c^{2}} \cdot \frac{\rho_{\mathrm{na}}}{b^{d} - 1}
  = \frac{\hbar}{c^{2}} \cdot \frac{\omega_{0}}{2\pi} \cdot 2\pi
  = \frac{\hbar\omega_{0}}{c^{2}}.
  \label{eq:massNA}
\end{equation}
\end{consequence}

\begin{proof}
Each full oscillation cycle of period $2\pi/\omega_{0}$ partitions the local
state into $b^{d}$ cells (by iterative refinement of Axiom~\ref{ax:bpsl}(iii)
over one period).  Exactly one cell is actualised per cycle --- the one the
system currently occupies; $b^{d} - 1$ are not.  The non-actualisation rate
is the number of non-actualised cells per unit time:
\[
  \rho_{\mathrm{na}}
  = \frac{b^{d} - 1}{2\pi/\omega_{0}}
  = (b^{d} - 1)\,\frac{\omega_{0}}{2\pi}.
\]
Dividing $\rho_{\mathrm{na}}$ by $(b^{d} - 1)/(2\pi)$ reduces it to
$\omega_{0}$; multiplication by $\hbar/c^{2}$ gives the rest mass
via~\eqref{eq:restmass}.  The result is~\eqref{eq:massNA}, which coincides
with Consequence~\ref{cons:restmass}.
\end{proof}

\begin{remark}[Interpretation]
Mass, in the deductive structure of this paper, is the normalised rate at
which a persistent bounded system fails to actualise all of its categorical
possibilities per unit oscillation time.  A more massive system accumulates
non-actualised cells faster; a lighter system accumulates them more slowly.
This is consistent with the $E = mc^{2}$ identity
(Consequence~\ref{cons:emc2}): rest energy is the information-theoretic cost
of the rate of accumulated distinctions.
\end{remark}

\subsection{Non-actualisation rate as a physical property}

An oscillator at rest frequency $\omega_{0}$ executes $\omega_{0}/(2\pi)$
cycles per second.  Each cycle partitions the local state space into $b^{d}$
categorical sub-cells; the oscillator occupies exactly one and vacates
$b^{d} - 1$.  The vacated cells are the ``non-actualised'' possibilities of
that cycle.  Over a macroscopic time $t$, the oscillator accumulates
$(\omega_{0}/2\pi)\cdot(b^{d} - 1)\,t$ non-actualised cells.

The oscillator's identity, in the formal sense, is the set of cells it has
\emph{not} actualised (Consequence~\ref{cons:loschmidt} and the associated
remark on identity from non-actualisation).  The rate at which this identity
accumulates is $\rho_{\mathrm{na}} = (b^{d} - 1)\,\omega_{0}/(2\pi)$ per
unit time.

This interpretation is consistent with the standard view of mass as
``resistance to acceleration'': an object with a higher non-actualisation
rate is harder to perturb, because each perturbation would require
re-actualising previously rejected possibilities, which costs energy by
Consequence~\ref{cons:compression}.  The mass is the coupling constant of
this energy cost to acceleration.

\section{Mass--Energy Equivalence}
\label{sec:emc2}

\begin{consequence}{Mass--Energy Equivalence}{emc2}
For a system with rest mass $m_{0}$ and momentum $p$, the relativistic
energy--momentum relation is
\begin{equation}
  E^{2} = (m_{0}c^{2})^{2} + (pc)^{2}, \qquad E_{0} = m_{0}c^{2}.
  \label{eq:emc2}
\end{equation}
\end{consequence}

\begin{proof}
From Consequence~\ref{cons:dispersion}, $\omega^{2} - c^{2}k^{2} = \omega_{0}^{2}$.
Multiply by $\hbar^{2}$:
\[
  (\hbar\omega)^{2} - c^{2}(\hbar k)^{2} = (\hbar\omega_{0})^{2}.
\]
By Consequence~\ref{cons:energyfrequency} ($E = \hbar\omega$) and its momentum
analogue ($p = \hbar k$),
\[
  E^{2} - c^{2}p^{2} = (\hbar\omega_{0})^{2} = (m_{0}c^{2})^{2},
\]
using~\eqref{eq:restmass}.  Rearranging gives~\eqref{eq:emc2}.  Setting
$p = 0$ (rest frame) gives $E_{0} = m_{0}c^{2}$.
\end{proof}

\begin{remark}
The formula $E = mc^{2}$ is usually presented as a consequence of special
relativity.  Here it is a Consequence of the Koopman ground-state energy being
non-zero on a bounded domain (Consequence~\ref{cons:restfreq}) combined with
the dispersion relation (Consequence~\ref{cons:dispersion}).  Special relativity
is not assumed; it is derived (Chapter~\ref{ch:lorentz}).
\end{remark}

\subsection{Decomposition of energy into rest and kinetic components}

Taking the square root of~\eqref{eq:emc2} and Taylor-expanding for small
momentum:
\begin{align*}
  E &= \sqrt{(m_{0}c^{2})^{2} + (pc)^{2}} \\
    &\approx m_{0}c^{2}\!\left(1 + \frac{p^{2}}{2m_{0}^{2}c^{2}}\right) \\
    &= m_{0}c^{2} + \frac{p^{2}}{2m_{0}}.
\end{align*}
The first term is the rest energy; the second is the non-relativistic kinetic
energy.  At high momentum ($p \gg m_{0}c$) the expansion breaks down and the
full relativistic form must be used.

\subsection{The photon as the massless limit}

Setting $m_{0} = 0$ in~\eqref{eq:emc2} gives $E = pc$; the rest energy
vanishes and all energy is kinetic.  Consequence~\ref{cons:dispersion} reduces
to $\omega = ck$, the dispersion relation of a massless mode propagating at
$c$.  The electromagnetic field, in our framework, is the massless Koopman mode
with $\omega_{0} = 0$; its quanta satisfy $E = \hbar\omega$ and
$p = \hbar k = E/c$.

\subsection{Binding energy reduces mass}

When a composite system has binding energy $E_{\mathrm{bind}} > 0$, its rest
mass is less than the sum of its constituents' rest masses by
$E_{\mathrm{bind}}/c^{2}$:
\[
  m_{\mathrm{composite}} = \sum_{i} m_{i} - \frac{E_{\mathrm{bind}}}{c^{2}}.
\]
This is the origin of the nuclear mass defect.  For example, $^{4}$He has
mass $4.002602$~u, less than four times the hydrogen atom mass
($4 \times 1.007825 = 4.031300$~u) by $0.028698$~u $= 28.3$~MeV/$c^{2}$,
matching the measured total binding energy \citep{audi2017}.

\section{Equivalence of Inertial and Gravitational Mass}
\label{sec:equivmass}

\begin{consequence}{Equivalence of Inertial and Gravitational Mass}{equivmass}
For every persistent bounded system, the inertial mass $m_{\mathrm{I}}$ and
the gravitational mass $m_{\mathrm{G}}$ are equal.
\end{consequence}

\begin{proof}
By Consequence~\ref{cons:restmass}, $m_{\mathrm{I}} = \hbar\omega_{0}/c^{2}$,
where $\omega_{0}$ is the rest frequency of the bounded oscillation.  The
gravitational mass is the coupling constant of the system to the ambient
partition depth field $\Depth(\mathbf{x})$ via the scalar coupling in the
partition Lagrangian.  The response of the system to a spatial gradient of
$\Depth$ is $\mu\ddot{\mathbf{x}} = -\nabla\Depth$, with $\mu = m_{\mathrm{I}}$
from the kinetic term of the partition Lagrangian.  Identifying the
gravitational force as the gradient response, $m_{\mathrm{G}} = m_{\mathrm{I}}$
identically.

The equivalence is a geometric identity: both masses are the same quantity
$\hbar\omega_{0}/c^{2}$, entering as the inertial parameter in the kinetic
term and as the coupling parameter in the scalar gradient term.  The
equivalence principle is not an independent postulate; it is forced by the
Lagrangian structure.
\end{proof}

\begin{remark}[Relation to general relativity]
Newton observed that inertial and gravitational mass were empirically equal
but gave no explanation.  General relativity elevates this observation to the
equivalence principle.  Consequence~\ref{cons:equivmass} derives the equivalence
from the partition Lagrangian: both masses are the same parameter $\mu$.  This
is the starting point for a derivation of geodesic motion in curved spacetime;
a full derivation would require identifying $\Depth$ (or a generalisation)
with the spacetime metric.  Such a derivation lies outside the scope of the
present work but is a natural continuation.
\end{remark}

\subsection{Weight and mass: inertial vs.\ gravitational}

Newton's distinction between weight ($W = mg$) and mass ($m$ in $F = ma$)
involves mass twice: once inertially, once gravitationally.  Newton observed
the equality empirically.  Consequence~\ref{cons:equivmass} explains why:
both roles are filled by the single parameter $\mu = \hbar\omega_{0}/c^{2}$
in the partition Lagrangian.  There is no separate ``gravitational charge'';
the same quantity that resists acceleration also couples to the depth
gradient.

\bigskip
\begin{tcolorbox}[colback=crownblue!5, colframe=crownblue!60!black,
                  title={\textbf{Chapter Summary: Mass as Accumulated Non-Actualisation}}]
\begin{itemize}[noitemsep]
  \item The rest mass of a bounded oscillator is $m_{0} = \hbar\omega_{0}/c^{2}$
    (Consequence~\ref{cons:restmass}), derived from the relativistic dispersion
    relation and the definition of group velocity.
  \item The Loschmidt principle (Consequence~\ref{cons:loschmidt}) establishes
    that the set of visited partition cells is monotone in time: entropy is
    non-decreasing.  The second law and the arrow of time are corollaries.
  \item Mass is the normalised rate of categorical non-actualisation per
    oscillation cycle: each cycle generates $b^{d} - 1$ unoccupied cells,
    and the normalised rate is $\hbar\omega_{0}/c^{2}$
    (Consequence~\ref{cons:massmemory}).
  \item Mass--energy equivalence $E^{2} = (m_{0}c^{2})^{2} + (pc)^{2}$
    follows from the dispersion relation by multiplying by $\hbar^{2}$
    (Consequence~\ref{cons:emc2}).
  \item Inertial and gravitational mass are the same parameter $\mu$ in the
    partition Lagrangian; the equivalence principle is derived, not postulated
    (Consequence~\ref{cons:equivmass}).
\end{itemize}
\end{tcolorbox}
